\section{Strengths and Limitations}

The first strength of this route is that it makes the service-pattern choice and the timetable artifact mutually interpretable. A weak timetable-first draft often draws a running diagram without making clear why that train structure exists. Here, the earlier passenger-flow and operation-plan layers explain the logic of the final schedule, so the timetable no longer appears as an isolated graphic but as the operational consequence of a demand-driven decision.

The second strength is that feasibility is treated as a separate evidence burden. Instead of compressing capacity, tracking, and dwell checks into one short assurance sentence, the route keeps these audits visible in the validation layer. This improves both credibility and paper rhythm: the result section answers what plan is recommended, while the validation section answers why that recommendation can still be trusted operationally.

The route also has limitations that should be stated explicitly. First, the service-pattern comparison depends on the way passenger demand is aggregated; if the overlap-zone burden is misestimated, then the large-small-route choice may already be distorted. Second, the timetable recurrence model usually simplifies real dispatching complexity by representing dwell, disturbance recovery, and rolling-stock circulation through a compact constraint system. Third, scenario conclusions remain conditional on the chosen perturbation range and on the assumption that the operating environment does not experience extreme disruptions.

For later improvement, two directions are especially important. One is to strengthen the comparison layer by adding richer candidate-plan baselines, so that the selected service structure is justified against more than one nearby alternative. The other is to extend the timetable layer toward rolling-stock circulation, disturbance absorption, or station-specific control policies. These extensions preserve the same paper chain while making the final recommendation closer to a competition-grade operational study.
